\appendix

\renewcommand\thesubsection{\thesection\arabic{subsection}}

\section{Appendix: Alternate Proofs of Exercises}

\subsection{Vector Spaces}

\begin{appendixexercise}[1c13]
    Prove that the union of three subspaces of $V$ is a subspace of $V$ if and only if one of the subspaces contains the other two.
\end{appendixexercise}

\begin{proof}[Coolempire93]
    Conversely, suppose that $U_1 \cup U_2 \cup U_3$ is a subspace of $V$ but none of the subspaces contains the other two. Let $W_1 = U_2 \cup U_3$ and $W_2 = U_1 \cup U_3$. By the assumption, we cannot have $W_1 \subseteq U_1$ or $W_2 \subseteq U_2$, so there exist $x \in U_1 \setminus W_1$ and $y \in U_2 \setminus W_2$. Now, since $U_1 \cup U_2 \cup U_3$ is a vector space, we have $x + \alpha y \in U_1 \cup U_2 \cup U_3$ for any nonzero $\alpha \in \mathbb{F}$. But if $x + \alpha y \in U_1$, then $(x + \alpha  y) - x = \alpha y \in U_1$, and as $\alpha$ is nonzero, this means that $y \in U_1$, a contradiction. Similarly, if $x + \alpha  y \in U_2$, then $(x + \alpha y) - \alpha y = x \in U_2$, a contradiction. Therefore, we must have $x + \alpha y \in U_3$ for all nonzero $\alpha$ to satisfy the closure of the vector space.

    As such, $x + y \in U_3$ and $x - y \in U_3$; that is, $(x + y) + (x - y) = 2x \in U_3$ and $(x + y) - (x - y) = 2y \in U_3$, contradicting the original definitions of $x$ and $y$$^\dagger$. Therefore, one of the subspaces must contain the other two.

    $^\dagger$ \textit{Final remark}. This is why the field must have more than two elements. Since part of the field axioms is that $0 \neq 1$, if the field only has two elements, it must be that $1 = -1$ (consider addition mod 2 as the addition operation), so $2x \in U_3$ would not mean that $x \in U_3$, as $2x = 0x = 0$.
\end{proof}